% Replacement for the current Section 4.1.
% This file is a LaTeX fragment, not a standalone document.

\subsection{The Black--Scholes Model as a Degenerate Measure-Ensemble Specialization}
\label{subsec:bs-degenerate-specialization}

The word \emph{degenerate} must be used relative to a precisely defined ambient
model.  In particular, the Black--Scholes model is not ``zero entropy'' at the
level of the stock price: even under its risk-neutral measure, the terminal
stock price is non-degenerate whenever $\sigma>0$ and $T>t$.  The collapse occurs
instead in the distribution over competing pricing measures and in the number
of active state-dependent pricing degrees of freedom.

\begin{definition}[Measure-ensemble pricing envelope]
\label{def:measure-ensemble-envelope}
Fix a filtered probability space
$(\Omega,\mathcal F,\mathbb F=(\mathcal F_t)_{0\leq t\leq T},P)$ and a strictly
positive numeraire
\[
  B_t=\exp\!\left(\int_0^t r_u\,du\right).
\]
Let $\mathfrak Q=\{Q^1,\ldots,Q^m\}$ be candidate measures equivalent to $P$,
let $w_i(t)\geq0$ be $\mathcal F_t$-measurable weights satisfying
$\sum_{i=1}^m w_i(t)=1$, and let $\Lambda_t$ be an $\mathcal F_t$-measurable
pricing-distortion operator.  For every payoff $H$ for which the following
conditional expectations exist, define
\begin{equation}
  \Pi^{\mathrm{ED}}_{t,T}(H)
  :=\Lambda_t\!\left(
       B_t\sum_{i=1}^m w_i(t)
       E^{Q^i}\!\left[B_T^{-1}H\mid\mathcal F_t\right]
     \right).
  \label{eq:ed-pricing-envelope}
\end{equation}
The associated meta-measure on the candidate-measure space is
\begin{equation}
  \nu_t:=\sum_{i=1}^m w_i(t)\,\delta_{Q^i}.
  \label{eq:meta-measure}
\end{equation}
Equation~\eqref{eq:ed-pricing-envelope} is a pricing envelope.  If it is to
define a dynamically consistent pricing process for arbitrary evolving
weights, additional tower-property or cocycle axioms must be imposed; those
axioms are not needed for the specialization result below.
\end{definition}

\begin{definition}[Black--Scholes degeneracy stratum]
\label{def:bs-stratum}
The Black--Scholes stratum of the envelope in
Definition~\ref{def:measure-ensemble-envelope} is specified by
\begin{equation}
  m=1,\qquad \nu_t=\delta_{Q^*},\qquad w_1(t)=1,
  \qquad \Lambda_t=\operatorname{Id},
  \qquad r_t\equiv r,\qquad \sigma_t\equiv\sigma>0.
  \label{eq:bs-stratum}
\end{equation}
Thus the $(m-1)$-dimensional simplex of narrative weights collapses to one of
its zero-dimensional vertices, the pricing distortion disappears, and the
coefficient processes are frozen at constants.
\end{definition}

\begin{theorem}[Black--Scholes specialization theorem]
\label{thm:bs-specialization}
Assume that $\mathbb F$ is the augmented filtration generated by a
$P$-Brownian motion $W^P$, and that a frictionless market contains the bank
account $B_t=e^{rt}$ and one stock satisfying
\begin{equation}
  dS_t=\mu S_t\,dt+\sigma S_t\,dW_t^P,
  \qquad S_0>0,\quad \sigma>0,
  \label{eq:physical-gbm}
\end{equation}
where $r,\mu,$ and $\sigma$ are constants.  Trading is continuous and
self-financing, and there are no portfolio constraints or transaction costs.
For any European payoff $H=h(S_T)$ satisfying
$B_T^{-1}H\in L^2(Q^*)$, the market
has a unique equivalent martingale measure $Q^*$ and its unique no-arbitrage
price is
\begin{equation}
  V_t
  =B_tE^{Q^*}\!\left[B_T^{-1}H\mid\mathcal F_t\right]
  =E^{Q^*}\!\left[e^{-r(T-t)}h(S_T)\mid\mathcal F_t\right].
  \label{eq:bs-risk-neutral-price}
\end{equation}
Moreover, equation~\eqref{eq:bs-risk-neutral-price} is exactly the restriction
of the measure-ensemble price~\eqref{eq:ed-pricing-envelope} to the degeneracy
stratum~\eqref{eq:bs-stratum}.  If $V_t=v(t,S_t)$ and $h$ has the standard
regularity (or is interpreted in the viscosity/weak sense), then $v$ satisfies
the Black--Scholes PDE
\begin{equation}
  \partial_t v+\frac12\sigma^2s^2\partial_{ss}v
  +rs\partial_sv-rv=0,
  \qquad v(T,s)=h(s).
  \label{eq:bs-pde-rigorous}
\end{equation}
Consequently, Black--Scholes is an exact lower-dimensional specialization of
the envelope, not merely an analogy to it.
\end{theorem}

\begin{proof}
Set
\begin{equation}
  \theta:=\frac{\mu-r}{\sigma},
  \qquad
  Z_t:=\exp\!\left(-\theta W_t^P-\frac12\theta^2t\right).
  \label{eq:bs-density}
\end{equation}
Because $\theta$ is constant, Novikov's condition holds and $(Z_t)$ is a
strictly positive $P$-martingale with $E^P[Z_T]=1$.  Define $Q^*$ by
\begin{equation}
  \left.\frac{dQ^*}{dP}\right|_{\mathcal F_T}=Z_T.
  \label{eq:qstar-definition}
\end{equation}
Girsanov's theorem implies that
\begin{equation}
  W_t^{Q^*}:=W_t^P+\theta t
  \label{eq:qstar-brownian}
\end{equation}
is a $Q^*$-Brownian motion.  Substituting
$dW_t^P=dW_t^{Q^*}-\theta\,dt$ into
equation~\eqref{eq:physical-gbm} gives
\begin{equation}
  dS_t=rS_t\,dt+\sigma S_t\,dW_t^{Q^*}.
  \label{eq:risk-neutral-gbm}
\end{equation}
It follows from It\^o's formula that
\begin{equation}
  d(B_t^{-1}S_t)=\sigma B_t^{-1}S_t\,dW_t^{Q^*},
  \label{eq:discounted-martingale}
\end{equation}
so $Q^*$ is an equivalent martingale measure.

It is unique.  Indeed, in a Brownian filtration the martingale representation
theorem represents the density process of any equivalent candidate martingale
measure by a stochastic exponential with some Girsanov kernel $\lambda_t$.
Under that measure the drift of $S$ is $\mu-\sigma\lambda_t$.  Requiring the
discounted stock to be a local martingale forces
\begin{equation}
  \mu-\sigma\lambda_t=r,
  \qquad\text{hence}\qquad
  \lambda_t=\frac{\mu-r}{\sigma}=\theta
  \quad(dt\otimes dP)\text{-a.e.}
  \label{eq:unique-kernel}
\end{equation}
Therefore every equivalent martingale measure has density
equation~\eqref{eq:bs-density}, and hence equals $Q^*$.  Equivalently, the one
non-degenerate risky asset spans the one Brownian source of risk; the martingale
representation theorem then yields replication of square-integrable claims.
This is the standard completeness--uniqueness relation for martingale models
\cite{harrisonpliska1983}.

By replication and no arbitrage, the value of $H$ must therefore be the unique
price in equation~\eqref{eq:bs-risk-neutral-price}.  Now impose
$\nu_t=\delta_{Q^*}$, $w_1(t)=1$, and
$\Lambda_t=\operatorname{Id}$ in
equation~\eqref{eq:ed-pricing-envelope}.  Direct substitution gives
\begin{align}
  \Pi^{\mathrm{ED}}_{t,T}(H)
  &=\operatorname{Id}\!\left(
      B_tE^{Q^*}[B_T^{-1}H\mid\mathcal F_t]
    \right)\\
  &=E^{Q^*}[e^{-r(T-t)}h(S_T)\mid\mathcal F_t]
   =V_t.
  \label{eq:exact-envelope-collapse}
\end{align}
This proves exact equality of the pricing operators on the degeneracy stratum.

Finally, equation~\eqref{eq:risk-neutral-gbm} is Markovian.  The Feynman--Kac
formula applied to equation~\eqref{eq:bs-risk-neutral-price} yields
equation~\eqref{eq:bs-pde-rigorous}; conversely, the discounted It\^o process
associated with a sufficiently regular solution of that PDE is a
$Q^*$-martingale and reproduces equation~\eqref{eq:bs-risk-neutral-price}.
This is precisely the Black--Scholes valuation equation
\cite{blackscholes1973}.  The claimed specialization follows.
\end{proof}

\begin{corollary}[Literal collapse limit]
\label{cor:bs-collapse-limit}
Let $Q^2,\ldots,Q^m$ be additional candidate measures and
\[
  \nu_t^\varepsilon
  =(1-\varepsilon)\delta_{Q^*}
   +\varepsilon\sum_{i=2}^m\alpha_i(t)\delta_{Q^i},
  \qquad \alpha_i(t)\geq0,
  \qquad \sum_{i=2}^m\alpha_i(t)=1.
\]
Define
\[
  X_t^\varepsilon:=B_t\int_{\mathfrak Q}
  E^Q[B_T^{-1}H\mid\mathcal F_t],\nu_t^\varepsilon(dQ).
\]
If all terms belong to a common normed pricing space and
$\|\Lambda_t^\varepsilon(X_t^\varepsilon)-X_t^\varepsilon\|\to0$ as
$\varepsilon\downarrow0$, then
\begin{equation}
  \Pi^{\mathrm{ED},\varepsilon}_{t,T}(H)
  \longrightarrow
  E^{Q^*}[e^{-r(T-t)}H\mid\mathcal F_t].
  \label{eq:literal-collapse-limit}
\end{equation}
Thus ``degeneration'' may be understood either as the exact boundary
restriction in Theorem~\ref{thm:bs-specialization} or as the explicit limit in
equation~\eqref{eq:literal-collapse-limit}.
\end{corollary}

\begin{remark}[Which entropy vanishes?]
\label{rem:meta-vs-state-entropy}
For the finite meta-measure~\eqref{eq:meta-measure}, define
\[
  \mathsf H_{\mathrm{meta}}(\nu_t)
  :=-\sum_{i=1}^m w_i(t)\log w_i(t),
  \qquad 0\log0:=0.
\]
On the Black--Scholes stratum,
$\nu_t=\delta_{Q^*}$ and hence
$\mathsf H_{\mathrm{meta}}(\nu_t)=0$.  This does not imply that the stock state
has zero entropy or zero variance.  In fact, conditional on $\mathcal F_t$,
\begin{equation}
  \log S_T
  \sim N\!\left(
      \log S_t+\left(r-\frac12\sigma^2\right)(T-t),
      \sigma^2(T-t)
    \right)
  \quad\text{under }Q^*,
  \label{eq:lognormal-nondegenerate}
\end{equation}
which is non-degenerate for $\sigma>0$ and $T>t$.
\end{remark}

\begin{remark}[Logical scope]
\label{rem:scope-of-bs-degeneracy}
The theorem proves an inclusion statement: once
equation~\eqref{eq:ed-pricing-envelope} is adopted as the ambient model,
Black--Scholes lies on its single-measure, frictionless, constant-coefficient
boundary.  It does not by itself prove that the larger envelope is uniquely
determined by the three emergence axioms, dynamically consistent away from
that boundary, empirically superior, or capable of producing arbitrage-free
prices for arbitrary choices of $w_i$ and $\Lambda_t$.  Those are separate
theorems.  Nor does the result make Black--Scholes ``false'': within its stated
assumptions it is the unique arbitrage-free valuation model for replicable
European claims.
\end{remark}

% Suggested bibliography entries/keys:
% blackscholes1973: F. Black and M. Scholes, The Pricing of Options and
% Corporate Liabilities, Journal of Political Economy 81(3), 1973, 637--654.
% harrisonpliska1983: J. M. Harrison and S. R. Pliska, A Stochastic Calculus
% Model of Continuous Trading: Complete Markets, Stochastic Processes and
% their Applications 15(3), 1983, 313--316.
